% Chapter 1 Turnitin revision blocks
% ----------------------------------
% The Chapter 1 source used to build the submitted PDF is not present in this
% workspace. Copy only the replacement body required at each marked location.
% Preserve the section/subsection labels already used by the main thesis file.

% REPLACEMENT 1
% Location: Chapter 1, Introduction, opening discussion (PDF thesis pp. 1--2;
% Turnitin report pp. 32--33). Replace the paragraphs beginning with
% "Healthcare is becoming increasingly dependent..." and ending with the
% description of the combined framework.

Healthcare institutions now generate longitudinal records that can support risk
prediction, monitoring, and clinical research. Evidence held by one hospital,
however, reflects that institution's patient population, measurement practice,
and care pathway. A model developed from one source can therefore lose accuracy
when it is applied elsewhere. Multi-institutional learning is valuable because it
can expose a model to a broader range of clinical variation, but conventional
data pooling creates additional privacy, ownership, consent, and breach risks.

Cross-silo federated learning changes the location of model fitting. Participating
institutions optimize locally and provide model contributions to a coordinating
process instead of transferring their source records \cite{mcmahan2017fedavg}.
Medical studies have demonstrated the feasibility of this arrangement across
natural institutional partitions, including brain-tumour segmentation and
COVID-19 outcome prediction \cite{sheller2020federated,dayan2021federated}.
Local retention nevertheless protects only one boundary. Shared parameters can
still expose information about training examples, and an authenticated
participant can still manipulate the collaborative model
\cite{shokri2017membership,zhu2019dlg,bagdasaryan2020backdoor}.

The engineering problem is consequently broader than deciding where data are
stored. A useful collaboration must limit the influence of individual patients,
conceal institutional updates from the coordinator, control the cost of encrypted
communication, detect or contain invalid state transitions, and preserve evidence
of who authorized a model release. This thesis investigates those requirements as
one bounded system and introduces CAABC-FL to connect adaptive compression with
the packing structure of CKKS ciphertexts \cite{cheon2017homomorphic}. The work
then evaluates predictive utility, privacy exposure, poisoning response,
communication cost, cryptographic correctness, and governance behaviour under a
common experimental protocol.

% REPLACEMENT 2
% Location: Section 1.1, Background of the Study, historical development
% (Turnitin report pp. 35--36). Replace the milestone paragraphs that describe
% FedAvg, membership inference, fully homomorphic encryption, CKKS, and secure
% aggregation.

Modern privacy-preserving collaboration emerged from several research streams
that developed independently. Fully homomorphic encryption established that a
program could, in principle, operate on encrypted values without first revealing
their plaintext \cite{gentry2009fully}. CKKS later supplied approximate packed
arithmetic suited to real-valued numerical workloads, making it relevant to model
aggregation \cite{cheon2017homomorphic}. Secure aggregation pursued a related
goal at protocol level: the coordinator learns a combined contribution while
individual client inputs remain hidden \cite{bonawitz2017practical}.

Federated optimization provided the learning structure into which these controls
could be inserted. FedAvg reduced the frequency of communication by allowing
clients to perform several local optimization steps before their contributions
were combined \cite{mcmahan2017fedavg}. Subsequent attacks demonstrated why
local data retention was not a complete privacy argument. Membership-inference
methods tested whether particular records influenced a trained model, while
gradient-inversion work showed that shared learning signals could reveal
approximations of private examples \cite{shokri2017membership,zhu2019dlg,
geiping2020inverting}. These milestones shifted the field from data-locality alone
towards layered protection of training, communication, aggregation, and release.

% REPLACEMENT 3
% Location: Subsection 1.1.1, State of the Art in the Modern Context, opening
% paragraphs (Turnitin report pp. 37--38).

Current research treats protected federated learning as a combination of
statistical, cryptographic, and operational controls. Differentially private
training limits the effect attributable to one protected unit; homomorphic or
secure aggregation conceals update values during coordination; and authenticated
services bind messages to recognized participants. In healthcare, multiparty
homomorphic analytics has already supported distributed survival and genomic
analyses, showing that protected computation can operate on substantive
biomedical tasks \cite{froelicher2021famh}. Reproducible frameworks and natural
multi-centre datasets have also improved the evaluation of institutional
heterogeneity \cite{ogier2022flamby,foley2022openfl}.

The remaining difficulty is compatibility. Many robust aggregators inspect the
distance, direction, or coordinates of every client update, whereas protected
aggregation deliberately hides those values. Encryption can also expand a small
logical update into several ciphertext objects, so a plaintext compression ratio
does not establish a network saving. Recent systems address parts of these
conflicts through encrypted anomaly detection, jointly generated privacy noise,
or privacy-compatible sparse aggregation \cite{xu2024ddfed,gu2025dpbrem,
asi2025preamble}. Their differing datasets, trust assumptions, and adversary
models leave room for an aligned healthcare evaluation that measures utility,
privacy, poisoning, encrypted bytes, and release governance together.

% REPLACEMENT 4
% Location: Section 1.2, Research Problem and Its Context (Turnitin report
% pp. 41--42). Replace the core problem statement.

Keeping records inside their originating hospitals does not make the complete
learning process private or trustworthy. An update can encode information about
the patients used to produce it; a malicious institution can submit a replacement,
sign-flipped, label-flipped, or backdoored contribution; and repeated transmission
of CKKS ciphertexts can impose considerably more traffic than the underlying
parameter vector. These risks interact. Concealing every individual update helps
confidentiality, for example, but prevents the coordinator from directly applying
many geometry-based poisoning defences.

The research problem is therefore to construct and evaluate a multi-institutional
learning framework that protects patient influence and institutional update
values, converts logical compression into a measurable encrypted-wire reduction,
and provides verifiable controls over participation, aggregation, and model
release. Existing work supplies strong components for individual parts of this
problem, but fewer studies test their compatibility under one dataset split,
threat model, and statistical protocol. The thesis addresses that integration gap
without assuming that differential privacy, encryption, or an audit record alone
is sufficient.

% REPLACEMENT 5
% Location: Section 1.4, Applications and Practical Use Cases, opening paragraph
% (Turnitin report pp. 45--46).

The proposed framework is intended as research infrastructure for an authorized
healthcare consortium rather than as an autonomous clinical device. Its immediate
use is to study how hospitals could build a common risk model while retaining
custody of source records and exchanging only protected, authenticated model
contributions. Potential tasks include early deterioration or sepsis prediction,
readmission analysis, multi-centre imaging research, and distributed clinical
analytics. Each application would still require task-specific validation,
institutional approval, and human oversight before operational use.

% REPLACEMENT 6
% Location: Section 1.6, Research Methodology Overview, first two paragraphs
% (Turnitin report pp. 49--51).

The study uses an artifact-centred experimental methodology. It first specifies
the actors, protected assets, adversary capabilities, and release conditions for a
cross-silo healthcare consortium. Those requirements are implemented as hospital,
coordinator, authority, verification, registry, and audit services. The learning
path combines patient-level private local optimization, CAABC-FL support
selection, CKKS aggregation, signed envelopes, independent aggregate
recomputation, validation, and rollback.

Evaluation uses the two natural source partitions of the 2019
PhysioNet/Computing in Cardiology Challenge sepsis cohort
\cite{reyna2020early,physionet2019sepsis}. Patients are separated before feature
construction so that observations from one person cannot enter both training and
evaluation. Comparisons are paired by seed and conducted under frozen protocols.
The analysis records clinical discrimination and calibration, serialized
ciphertext traffic, privacy-accounting output, registered attack outcomes,
cryptographic error, latency, and governance failure cases. Confidence intervals
are calculated over independent model seeds rather than treating predictions or
attack trials from one trained model as independent experiments.

% REPLACEMENT 7
% Location: Section 1.7, Research Contribution, introductory paragraph
% (Turnitin report pp. 51--53). Keep the existing detailed contribution list
% after this paragraph.

The contribution of this research lies in connecting mechanisms that are usually
evaluated separately and then testing the resulting trade-offs. The work provides
a formalized end-to-end architecture, the CAABC-FL communication mechanism, an
executable multi-service research system, and a paired experimental comparison on
natural hospital partitions. Its central empirical claim is deliberately narrow:
under the registered protocol, CAABC-FL reduces serialized encrypted upload while
remaining within the specified AUROC non-inferiority margin. Privacy, AUPRC,
poisoning, computation, and governance outcomes are retained as separate evidence
and are not used to claim universal superiority.

% REPLACEMENT 8
% Location: Sections 1.8, 1.8.1 and 1.8.2 (Turnitin report pp. 53--55).

The project was completed with existing computing resources and openly available
research software. Consequently, the budget records the resources actually used
rather than assigning an artificial monetary value to free tools or public data.
No dedicated server, paid clinical dataset, cloud subscription, laboratory
equipment, or project travel was required. This cost statement applies only to
the reported research prototype; an independent multi-hospital deployment would
require a separate budget for infrastructure, security operations, governance,
and clinical validation.

The timeline groups the work into problem definition, literature review, data
preparation, system design, implementation, controlled experimentation,
statistical analysis, and thesis preparation. Some activities overlap because
experimental findings were used to refine the threat model and implementation.
The Gantt chart records that iterative progression and distinguishes development
work from the later confirmatory evaluation.

% REPLACEMENT 9 -- CRITICAL CONTENT CORRECTION
% Location: Section 1.9, Organization of the Dissertation (Turnitin report
% p. 56). Replace the entire existing section because it still describes a heart
% disease prediction thesis.

The remainder of the dissertation is organized into four chapters. Chapter~2
reviews cross-silo federated learning, differential privacy, homomorphic and
secure aggregation, communication-efficient optimization, privacy attacks,
poisoning defences, and accountable governance. It then compares representative
studies and identifies the integration and evaluation gaps addressed by this
research.

Chapter~3 translates those gaps into the conceptual framework, threat model,
problem formulation, layered architecture, CAABC-FL algorithm, data-processing
pipeline, and experimental protocol. It also states the assumptions that separate
the implemented research system from a production clinical deployment.

Chapter~4 reports predictive, communication, privacy, poisoning, cryptographic,
and service-level results. It presents paired statistical analyses, same-protocol
baseline comparisons, ablation evidence, negative outcomes, and the limits of the
supported claims. Chapter~5 consolidates the findings and contributions, records
the study's limitations, and identifies the technical and clinical work required
for stronger future deployment evidence. Supplementary protocols, measurements,
and engineering mappings are provided in the appendices and annexures.
